% Struss 2008 Chapter 10; source-specific construction, not Chapter 9 ATMS formulation.
\chapter{Compositional Models and Model-Based Problem Solving}
\label{sec:p16-qualitative-compositional-model-formulation}

Peter Struss's account starts from a separation between reusable domain
knowledge and the reasoning which uses it for a task
\cite[secs.~10.1--10.3]{Struss2008ModelBasedProblemSolving}.
A library describes how constituents behave; a system description says how
they are connected; diagnosis, testing, design and remedy engines exploit
the resulting models differently. This is not a requirement to express every
physical model in first-order logic. Consistency and entailment are semantic
requirements which a constraint solver, numerical predictor or other modeling
formalism must implement at the needed scope.

\section{The architecture and the task determine what a model must do}
\label{sec:p16-cm-tasks}

In the source's vehicle example, a rule excluding a flat battery whenever
either set of lights is on works only because starter and lights share a
supply. Figure~10.1 changes the circuit to two batteries: illuminated lights
can then coexist with a flat starter battery. The defect is the compiled-in
structure, not the rule syntax. A model library for batteries, wires, lights
and starters can instead be instantiated against either connection structure.
The same library can support prediction of failure effects or test generation,
not only diagnosis \cite[pp.~395--397]{Struss2008ModelBasedProblemSolving}.

Let \(M\) denote a behavior model, \(O\) observations and \(G\) goals.
The following task distinctions are consequential:
\begin{itemize}
\item Diagnosis searches for a revised model consistent with \(O\), or,
      under a stronger explanatory requirement, entailing selected symptoms.
      Reducing a goal-violation assessment to checking only the all-correct
      model presumes that all-correct behavior guarantees the goals.
\item Design seeks \(M\entails G\). Consistency of \(M\cup G\) is
      useful for refuting alternatives, but only says a desired behavior is
      possible. Configuration fixes much of the structure and selects types
      or parameters; unrestricted innovative design can change the structure.
\item Failure analysis asks whether a fault model entails an unwanted effect
      or merely permits it in a specified scenario. Possible is not inevitable.
\item A test chooses a stimulus and observations which separate remaining
      hypotheses. A measurement proposal holds the stimulus fixed.
\item A remedy chooses actions or a target configuration consistent with,
      or guaranteed to achieve, goals. It may need to work for several
      still-possible diagnoses and respect intermediate safety goals.
\end{itemize}
Figure~10.2's shared architecture is a predictor/consistency checker coupled
to a model-revision generator. Disturbances, design alternatives and remedial
interventions supply different revision spaces, not interchangeable hypotheses
\cite[sec.~10.2]{Struss2008ModelBasedProblemSolving}.

\section{Reusable fragments and actual model composition}
\label{sec:p16-cm-composition}

Struss distinguishes six modeling requirements: prediction/consistency
checking, behavioral validity, explicit conceptual structure, automated
composition, genericity, and appropriate granularity. A global equation
without its component/process provenance cannot say which physical hypothesis
caused an inconsistency. A predictor fixed to one input-to-output direction
can miss contradictions requiring reasoning from measured outputs.
An overly restrictive model can falsely reject real behavior; a conservative
model may keep additional possibilities but should not exclude the actual one
\cite[secs.~10.3.1--10.3.6]{Struss2008ModelBasedProblemSolving}.

\begin{hbdefinition}{def:p16-cm-library}{Component modes, guarded fragments and connections}
Let \(C\) be a finite component set. Each component \(c\) has a finite,
nonempty mode set \(B_c\), containing \(ok_c\). Exactly one mode holds
for each component. A complete mode assignment \(a\in\prod_{c\in C}B_c\)
selects one mode per component; a partial assignment selects modes on a subset.
Assignments also denote the conjunctions of their selected mode atoms.
The library \(L\) contains variable-domain axioms, mode exclusivity and
coverage, and guarded fragments \(b_c\to F_{c,b}\). The system structure
\(S\) supplies connections, typically equality of identified port variables.
The composition for \(a\) is \(L\cup S\cup\{a\}\).

In a relational realization, mode \(b\) supplies a relation \(R_{c,b}\)
on local variables. Embed it in the global variable domain \(D\), enforcing
the connections, and take their natural join/intersection \(R_a\).
Retaining a port equality is different from independently choosing two
values which happen to be in the same qualitative category.
\end{hbdefinition}

The library is domain-specific; the connections are system-specific.
Struss's \(SD=LIB\cup STRUCTURE\) abbreviates them but does not erase
the distinction. Mode exclusivity does not prove that every composed model
is consistent. That additionally needs compatible domains, connections and
valid fragments. The formal \(a\) above is an assignment of conceptual
behavior modes, not an arbitrary fit of numerical parameters.

\begin{hbtheorem}{thm:p16-cm-composition}{proved=annex:p16.cm-composition}{Conservative composition and exact observable compilation}
Suppose each actual global behavior in an application class satisfies the
connections and projects into the selected local mode relations. Then it
belongs to \(R_a\); a genuine actual observation cannot be refuted by this
composition. For observable variables \(V\), the exact compiled model
\(P_a=\pi_V(R_a)\) preserves consistency with every restriction on \(V\)
and every universal prediction expressed only on \(V\), in both directions.
Independent projection of each fragment before joining need not do so.
\end{hbtheorem}

For example, with binary hidden \(z\), one fragment requires \(z=x\)
and another \(z=y\). Their join entails \(x=y\). Separately deleting
\(z\) leaves both observable domains unrestricted and admits \(x=0,y=1\).
Eliminating hidden variables from the composite model is thus a compilation
operation, not permission to discard inter-fragment coupling.

\begin{hbexample}{def:p16-cm-genericity}{A pipe, a motor and query-relative granularity}
A pipe beside a check valve must still represent reverse flow: a broken
check valve can expose the pipe to that environment. A tram motor can work
as a generator during braking. Encoding the intended global direction in
the local fragment defeats reuse. Relevant exogenous values must be covered
over their declared domains, not only those seen during normal operation.

Fine structural resolution is needed to distinguish the faults which matter
to repair; fine numerical resolution is needed only when it changes the
task's distinctions. For fixed on-board sensors, compilation can eliminate
many internal variables. For workshop diagnosis with movable probes, that
same compilation may discard variables later needed. Struss's pipe-library
challenge separates oxygen transport in an intake, emissions in an exhaust,
and pressure supply to a valve: one overloaded fragment is not automatically
an adequate tailored model \cite[pp.~406--407,459]{Struss2008ModelBasedProblemSolving}.
\end{hbexample}

\section{Mode search, conflicts, fault models and relevance}
\label{sec:p16-cm-mode-search}

A complete \(a\) is a consistency diagnosis when \(L\cup S\cup\{a\}
\cup O\) has a model. The faulty-component set is
\(\{c:a(c)\ne ok_c\}\). Fault localization abstracts from the particular
fault modes; fault identification retains them. Inclusion-minimal faulty sets,
minimum cardinality and most probable assignments are different criteria.
The normality-only conflict/hitting-set construction and its corrected
shared-node implementation are developed in
\cref{sec:p16-consistency-based-diagnosis}. Here the important additional
case is a library with actual fault modes
\cite[secs.~10.4.1--10.4.2]{Struss2008ModelBasedProblemSolving}.

\begin{hbdefinition}{def:p16-cm-conflict}{A mode conflict and its use}
An inconsistent partial assignment \(p\) satisfies
\(L\cup S\cup O\cup\{p\}\entails\bot\). Its conflict is the
clause \(\neg p\): at least one of its selected modes must change.
A smaller inconsistent assignment gives a stronger pruning clause. Minimal
conflicts are useful, but a sound nonminimal one is sufficient for safe pruning.
\end{hbdefinition}

\begin{hbtheorem}{thm:p16-cm-mode-conflicts}{proved=annex:p16.cm-mode-conflicts}{Conflict characterization and the fault-model boundary}
Let \(\mathcal K\) consist of the negations of all inclusion-minimal
inconsistent partial mode assignments, including the empty one if the
background is inconsistent. A complete assignment is a diagnosis exactly
when it satisfies every clause in \(\mathcal K\). A partial assignment
is a kernel exactly when it is inclusion-minimal among terms whose every
complete mode extension satisfies \(\mathcal K\). This is implication
relative to the mode-domain axioms, not arbitrary unconstrained Boolean
valuations.

In the two-mode specialization \(ok_c,fault_c\), the family of feasible
faulty-component sets is upward closed exactly when every minimal
inconsistent partial assignment contains only \(ok\) assumptions.
Nonminimal redundant conflicts need not have that restricted syntax.
With actual fault-mode constraints, neither upward closure nor generation
of all diagnoses by minimal faulty sets is automatic.
\end{hbtheorem}

This states the content behind the source's basic-conflict and kernel
characterizations at a precise finite mode scope. Its positive-conflict
hitting-set account applies when the relevant minimal conflicts contain
only normality assumptions; the complete normality-only proof is in
\cref{thm:p16-diag-conflicts}. Fault-mode constraints must not be silently
removed to reuse that simpler search. Struss credits de Kleer, Mackworth
and Reiter for these general characterizations, and Reiter for the
normality-based hitting-set account
\cite[Theorems~10.1--10.3 and references~19,62]{Struss2008ModelBasedProblemSolving}.

\begin{hbmethod}{def:p16-cm-candidate-search}{Constructing and checking candidates}
Generate complete mode assignments in an order appropriate to the task.
Before instantiating a candidate, test cached conflict clauses. If it extends
an inconsistent partial assignment, discard it. Otherwise compose its mode
fragments with the structure and observations and call the consistency checker.
Record a consistent candidate, or extract/cache a supporting inconsistent
partial assignment. Revise modes rather than changing unrelated mathematical
constraints without conceptual provenance. After new observations, recheck
survivors and also check every newly generated candidate against all old data.
\end{hbmethod}

\begin{hbtheorem}{thm:p16-cm-search}{proved=annex:p16.cm-search}{The finite candidate construction with explicit oracle duties}
For \(N=\prod_{c\in C}|B_c|\) candidates, a fair enumeration without
repetition and a total exact consistency oracle terminates after at most
\(N\) oracle calls, returns exactly the consistency diagnoses, and remains
so when candidates extending sound cached conflicts are skipped.
Filtering the returned set by any partial order on this finite family yields exactly
its maximal candidates. A linear order extending the strict preference lets
one skip candidates already dominated by a recorded diagnosis when only
maximal diagnoses are wanted. This last skipping requires that order.
These are finite-space guarantees, not an efficiency theorem for DDE,
SHERLOCK, arbitrary first-order consistency or unbounded process creation.
\end{hbtheorem}

Propagation plus dependency recording, such as an ATMS, can find conflicts
from incompatible predicted values. In resistive or hydraulic models,
propagation may be incomplete. Missing a conflict keeps extra candidates;
soundness of discovered conflicts alone does not certify consistency of a
survivor. Compilation can improve runtime, but may shift prohibitive work
offline. Hierarchical diagnosis needs valid aggregate normal and fault models;
hand-building them forfeits much of the intended reuse
\cite[pp.~420--422]{Struss2008ModelBasedProblemSolving}.

\begin{hbexample}{def:p16-cm-lights}{The source's multiple-fault circuit and a fault-model revision}
Figure~10.3 has a dark rear light \(R\), inactive starter \(T\), and
lit head light \(H\), with battery \(B\) and wires \(w_1,\ldots,w_6\).
The four reported normality conflicts have component supports
\[
\begin{split}
K_1&=\{B,w_1,w_2,T\},\\
K_2&=\{B,w_1,w_2,w_3,w_4,R\},\\
K_3&=\{H,w_5,w_6,R\},\\
K_4&=\{H,w_5,w_6,w_3,w_4,T\}.
\end{split}
\]
The pair \(\{R,T\}\) hits all four; neither singleton does.
Deriving \(K_3,K_4\) uses the lit head light to infer supply information,
not only forward causal propagation. Normality-only models also permit
implausible explanations involving a defective lit head light.

If the library exhaustively says every nonnormal head-light mode is dark,
its observed illumination refutes them and exonerates that component.
But an unrestricted unknown fault cannot be refuted this way. The source's
exotic open/short-to-source wire fault demonstrates why an incomplete catalog
must not be treated as physical exhaustiveness. With explicit fault models,
adding \(H\) to \(\{R,T\}\) can make the candidate inconsistent:
upward closure of faulty sets is no longer available
\cite[pp.~412--415]{Struss2008ModelBasedProblemSolving}.
\end{hbexample}

\begin{hbdefinition}{def:p16-cm-circuit-realization}{An explicit voltage-presence realization of the source circuit}
To make that calculation reproducible, give the battery and three branches
binary left/right potentials \((x_j,y_j)\), \(j=B,T,R,H\).
Normal wires identify potentials on the two rails: \(w_1,w_2\) connect
\(B\) to \(T\), \(w_3,w_4\) connect \(T\) to \(R\), and
\(w_5,w_6\) connect \(R\) to \(H\), with the odd wire on the left.
Normal battery behavior requires \(x_B\ne y_B\). A normal starter/light
is active/lit exactly when its two potentials differ. Observations are
starter inactive, rear light dark and head light lit. A faulty component
adds no constraint; the port domains and observations still hold.
This is an explicit ideal voltage-presence model for the source topology,
not a quantitative circuit or a catalog of physically exhaustive fault modes.
\end{hbdefinition}

\begin{hbtheorem}{thm:p16-cm-circuit}{proved=annex:p16.cm-circuit}{The four source conflicts and every minimal localization}
For that model, the minimal normality conflicts are exactly \(K_1\)
through \(K_4\). Put \(A_0=\{B,w_1,w_2\}\),
\(C_0=\{H,w_5,w_6\}\), and \(D_0=\{w_3,w_4\}\).
All minimal fault localizations are
\[
 \{R,T\},\qquad
 \{a,c\}\ (a\in A_0,c\in C_0),\qquad
 \{a,R,d\}\ (a\in A_0,d\in D_0),\qquad
 \{c,T,d\}\ (c\in C_0,d\in D_0).
\]
There are \(1+9+6+6=22\): the displayed \(\{R,T\}\) and
twenty-one alternatives. The count is of inclusion-minimal localizations
under unrestricted faulty behavior, not of physically plausible failures
or probability-ranked mode assignments.
\end{hbtheorem}

An exhaustive physical-negation axiom says
\(\neg ok_c\to\bigvee_b fault_{c,b}\). Its exoneration inference
requires all those alternatives to be refuted in the same context. It is
not justified by refuting only the likely ones. A kernel diagnosis is a
minimal partial mode assignment all of whose complete extensions are
diagnoses; it is not an ordinary minimal faulty-component set.

\begin{hbtheorem}{thm:p16-cm-equal-chain}{proved=annex:p16.cm-equal-chain}{The source's chain and the size of its generator}
In a chain of \(n=17\) components with increments \(0,-1,+1\),
input zero and final output one, a faulty-component set is feasible exactly
when its cardinality \(k\) is odd. For that set the complete diagnoses have
\((k+1)/2\) positive and \((k-1)/2\) negative increments. There are
\(17\) singleton diagnoses but \(2^{16}\) feasible faulty sets overall.
Every kernel assignment must specify all components. Thus all complete
diagnoses are needed as kernels, with multiplicity of mode assignments for
larger faulty sets; fault localizations and kernels are not identical objects.
\end{hbtheorem}

\section{Preferences rather than an arbitrary consistent model}
\label{sec:p16-cm-preferences}

The source's mode preference is a partial order for each component, with
\(ok\) highest and an unrestricted unknown mode, if present, lowest.
The product order compares complete assignments coordinatewise. It does
not rank incomparable faults in different components. DDE explores successors
only after preferred alternatives have been refuted; SHERLOCK adds a global
probability scale and assumes independence when forming mode-assignment
priors. A stopping threshold tied to posterior mass can return only selected
candidates, not the exhaustive answer \cite[pp.~415--420]{Struss2008ModelBasedProblemSolving}.

\begin{hbtheorem}{thm:p16-cm-default-sound}{proved=annex:p16.cm-default-sound}{The source preference-default implication at finite mode scope}
Let \(W=L\cup S\cup O\) be consistent, with the finite exclusive
mode families above. For every mode \(b\) form the normal default with
prerequisite \(\bigwedge_{d>b}\neg d\), justification \(b\), and
conclusion \(b\). Empty prerequisite means truth. Use Reiter's extension
semantics, including least prerequisite-driven generation, not arbitrary
self-support. If \(\operatorname{Th}(W\cup\{a\})\) is an extension,
then the complete assignment \(a\) is a product-preferred diagnosis.
This is the implication in Struss's Theorem~10.4; no converse or claim that
local defaults rank all globally probable diagnoses is asserted.
\end{hbtheorem}

The preference-default construction is credited to Dressler and Struss's
1992 work in the chapter; DDE's later control strategies and the probability
search in SHERLOCK are separate implementation referrals. This proof is
not a claim to have read or reproduced those entire original publications.

The source's conflict-directed search is more specific than calling a SAT
solver once: the task usually needs a relevant set of explanations, support
for new observations, and sometimes explicit exclusion of safety-critical
faults. Relevance is governed by the downstream action and risk, not by a
universal numerical confidence attached to logical consistency.

\section{Diagnosis across time and transition composition}
\label{sec:p16-cm-temporal}

A temporal hypothesis is a history of complete mode assignments on an
observation interval, not one assignment independently chosen at each sample.
The observation history \(O_H\) records values and their times. The model
also constrains possible state sequences and mode changes; boundaries of
the history intervals require an explicit transition convention. A temporal
diagnosis has a full history satisfying the model and \(O_H\).
A state-based diagnosis instead checks the relevant mode assignment against
each observed snapshot \cite[pp.~423--427]{Struss2008ModelBasedProblemSolving}.

\begin{hbtheorem}{thm:p16-cm-temporal}{proved=annex:p16.cm-temporal}{Snapshot necessity and the missing inter-snapshot condition}
A temporal diagnosis restricts to consistent observed snapshots. The converse
does not follow just from a persistent mode and homogeneous decomposition
into state constraints and a common set of sequence constraints. For example,
the persistent model \(x'=0\) admits each snapshot \(x(0)=0\),
\(x(1)=1\) separately, with instantaneous derivative zero, but no
differentiable trajectory satisfying both. Sequence constraints or their
observation consequences must still be enforced across the samples.
\end{hbtheorem}

In the source's parked-car case, both stationary endpoint observations can
fit an intact parking brake, but different positions imply an intermediate
nonzero speed. For continuously differentiable position, the mean value
theorem supplies that intermediate speed. This indirect observation refutes
the intact-brake model without separately simulating every fault candidate.
It does not identify the departure time or trajectory.

Struss distinguishes persistence of a mode from a time-independent input/output
function. A persistent fault remains after occurrence unless repaired; it
may generate changing output. Homogeneity means all mode models share the
same sequence constraints. Discrete OPEN/CLOSED transitions allowed in a
normal valve but disallowed by STUCK violate that commonality. The source's
state-based sufficiency discussion refers to additional observation coverage
and sequence-derived information; it must not be read as a universal claim
that homogeneity alone glues arbitrary snapshot witnesses.

For persistent mode candidates, extend observations by sound consequences
of shared sequence constraints when possible, then check candidate state
constraints against that extended data. New conflicts can be accumulated.
Newly generated candidates must also face previous observations, not just
the newest sample. Equality between this construction and full temporal
diagnosis requires the appropriate observation/model conditions; merely
adding finitely many estimated derivatives is not a general proof of it.

\begin{hbdefinition}{def:p16-cm-transition-model}{Component transition models and a finite trace specialization}
The source gives a component model with finite states \(S_c\), initial
state \(s_{c,0}\), events \(E_c\), observable events, transitions
\(T_c\subseteq S_c\times E_c\times S_c\times E_c\), and fault
transitions. Trigger and emitted event are distinct fields. System states
are products of local states; aggregate events can be sets of local events.
Synchronous versus asynchronous composition is a modeling decision.

For the finite trace construction below, fix that decision and its resulting
finite graph \(Q\), initial state \(q_0\), and finite fault-label set
\(F\). Label every edge by one observable symbol or \(\epsilon\),
and by the subset of fault labels occurring on that edge. This is a
specialization after the event convention is fixed, not a hidden assumption
that every source transducer emits exactly one observation per transition.
An observed word must be the complete observable projection of a path rooted
at \(q_0\), with no unexplained additional observed symbols.
\end{hbdefinition}

\begin{hbmethod}{def:p16-cm-trace-construction}{The finite hidden-closure observer}
Keep pairs \((q,J)\), where \(q\) is a graph state and \(J\subseteq F\)
the faults seen so far. Start from \((q_0,\varnothing)\) and close under
\(\epsilon\)-edges, unioning their fault labels. For each next observed
symbol, take all matching labeled edges, update the fault sets, and again
close under hidden edges. A closure adds a pair only if it is new.
The resulting set describes both possible current states and possible
fault sets. It intentionally forgets multiplicity and order of repeated
fault events; a task requiring those needs a richer history representation.
\end{hbmethod}

\begin{hbtheorem}{thm:p16-cm-trace}{proved=annex:p16.cm-trace}{Exact finite-path coverage of the stated observer}
After any finite observed word, the observer contains exactly the end-state
and accumulated-fault-set pairs of finite paths with that observable
projection. Hidden closure terminates after at most \(|Q|2^{|F|}\)
pair insertions. Fault detection is justified when the resulting set is
nonempty and every pair has nonempty fault set. An empty set refutes the
model/observation combination; it is not positive evidence for a fault.
\end{hbtheorem}

Figure~10.4 illustrates why hidden faults must remain represented. After
\(\alpha\), the diagnoser contains states \(3\) and \(7\) labeled
faulty together with state \(11\) labeled normal. A second \(\alpha\)
leaves state \(6\), labeled faulty. The first observation is ambiguous;
the second detects the fault. Storing only the most likely path would not
have this exhaustive trace interpretation. Struss attributes this diagnoser
figure to Sampath, Lafortune and Teneketzis; the transition formalism is not
the probability-tracking policy of Livingstone.

The discrete-event diagnoser compiles observation-language behavior.
Livingstone instead composes shared physical/control variables and synchronizes
component transitions, then incrementally tracks selected likely paths.
Its qualitative state constraints turn continuous distinctions into finite
ones, but probabilities, transition semantics and abstraction validity remain
additional ingredients. Local finite component sets do not prevent the
Cartesian-product explosion \cite[pp.~428--431]{Struss2008ModelBasedProblemSolving}.

\section{Abduction and process-oriented situation assessment}
\label{sec:p16-cm-process-assessment}

Struss's abductive diagnosis separates conditioning observations \(O_C\)
from observations to be explained \(O_A\). It requires
\(M\cup O_C\) to be consistent and to entail \(O_A\).
Commands are not automatically symptoms which a fault should explain.
Logical entailment by itself is not a causal relation
\cite[sec.~10.4.5]{Struss2008ModelBasedProblemSolving}.

\begin{hbexample}{def:p16-cm-parallel-bulbs}{The missing supply premise in Figure 10.5}
Condition on a CLOSE command for a switch feeding two parallel bulbs.
With bulb one broken and bulb two normal, the model makes the first dark
but leaves the second's state dependent on the unknown supply. In a binary
idealization with supply \(u\), outputs are \((0,u)\): both \(u=0\)
and \(u=1\) satisfy the fault and command. Thus the observations
\((OFF,ON)\) are consistent, witnessed by \(u=1\), but not entailed,
as \(u=0\) is a countermodel. Adding/abducing the supply state \(u=1\)
produces the explanatory entailment. Merely relabeling the same fault
assignment abductive does not supply that premise.

The source also shows representation sensitivity: a model predicting
zero-or-positive flow does not entail strictly positive flow. Coarsening both
observation and model to nonnegative makes the entailment available, but
loses the zero/positive distinction needed by other tasks. That is a change
in the explained claim, not proof of the stronger original symptom.
\end{hbexample}

\begin{hbdefinition}{def:p16-cm-process-fragment}{Processes, influence closure and introducible boundaries}
A process fragment specifies typed participating objects, structural conditions,
quantity conditions, structural effects and quantity effects. Matching types
and relations creates candidate instances; satisfying both condition sets
activates an instance. Structural effects can create objects and relations.
Quantity effects are contributions, not necessarily the final net direction.
In Struss's notation \(I^+(x,y)\) means \(x\) positively influences
\(y'\): \(y'=f(\ldots,x,\ldots)\) with positive partial derivative.
This orders the arguments differently from the affected-quantity/rate notation
used in the native QPT heat-flow example. Combining all contributions requires
an explicit closed-world assumption about the influence set.

An introducible object type can occur at the explanatory boundary without
being generated by another modeled process. Other new objects need a
generating process and its conditions. The library and task determine that
boundary; an arbitrary unexplained object at an internal point is not allowed.
\end{hbdefinition}

\begin{hbmethod}{def:p16-cm-process-revision}{The source's structural and assumption revision workflow}
Start from observed objects/values and defeasible scenario assumptions.
Instantiate matching process fragments and resolve influences using recorded
closure assumptions. If the model contradicts observations, extract a
conflict among user assumptions and/or influence closures. Retract a relevant
assumption. Retracting an influence closure requires searching the library
for a process capable of contributing to that typed quantity. Satisfy its
quantity conditions and structural conditions. Missing structural participants
must be permitted introducibles or generated by further process instances.
Recompose and check the new model; a new inconsistency starts another revision.

Prefer assessments which retract no unnecessary assumptions and introduce
no unnecessary boundary objects. These are two dimensions of inclusion
minimality, not a guaranteed one-step heuristic based on the number of
newly instantiated processes. Relations are not counted in the source's
particular minimality criterion. A finite process-type library does not bound
the number of object instances or prove termination of this revision search.
\end{hbmethod}

\begin{hbexample}{def:p16-cm-iron}{The source-native iron explanation and its explanatory boundary}
Figure~10.6 follows reservoir water through a pump, tank and treatment plant
to the tap. The observation is an unwanted high dissolved-iron concentration.
The assessment introduces solid iron in the sediment, retracts the usual
neutral-pH assumption for the lower layer, and activates iron redissolving
under acidic conditions. Dissolved iron is generated by that process, then
transported to upper layers, intake, tank and tap. Solid sediment iron is a
boundary object; dissolved iron in the tank needs the process/transport account.
The source suggests algal bloom as a possible earlier source of acidification,
not as an independently established observation in every assessment.

Redissolving contributes positively to dissolved iron, but oxidation could
counteract it. Its presence alone does not prove a positive net derivative.
Nor can the high concentration be explained merely by assigning a pump a
generic broken mode. The environmental part is not a well-designed artifact:
the drinking-water goal is external, and natural processes can violate it
without being faulty components. Figure~10.7 makes assumptions and influence
closures, not just component modes, explicit revision targets
\cite[pp.~434--437]{Struss2008ModelBasedProblemSolving}.
\end{hbexample}

The situation-assessment triple consists of structure, quantity assignments
and retracted assumptions. It must be consistent with observations and retained
assumptions, retain observed objects, explain nonboundary objects by process
effects, and satisfy the stated minimality criterion. This is a generative
structural search, not the fixed Cartesian product of component modes.
The chapter's PDE and G+DE descriptions are snapshot methods: an earlier
algal bloom can cease while its iron consequences persist. A present-time
absence must not be used to rule out that historical mechanism without a
temporal model \cite[pp.~436--438]{Struss2008ModelBasedProblemSolving}.

Control-engineering fault isolation is another distinct regime. The source
describes numerical residual/model revisions, often parameter changes in a
fixed black-box structure, without the same explicit conceptual component/
process layer. Such models can be appropriate to their applications, but
their parameter search is not a replacement for structural fault hypotheses
or newly introduced substances \cite[sec.~10.4.7]{Struss2008ModelBasedProblemSolving}.

\section{Constructing discriminating tests from the composed relations}
\label{sec:p16-cm-tests}

For testing, a candidate hypothesis is a relation \(R_i\) on system
variables. With component hypotheses it is the join of the selected local
mode relations. Observable projection removes variables which cannot be
measured. The stimulus variables are included among the observables because
the applied stimulus is known. Figure~10.8 depicts intersections of observable
behavior relations, not geometric overlap of physical components
\cite[pp.~439--443]{Struss2008ModelBasedProblemSolving}.

\begin{hbdefinition}{def:p16-cm-test-relations}{Admissible stimuli and observable response sets}
Let \(U\) be the stimulus domain and \(Z\) the remaining observable
response domain. Write \(P_i=\pi_{U,Z}(R_i)\) and
\(P_i(u)=\{z:(u,z)\in P_i\}\). Restrict to admissible stimuli
\(A\subseteq U\), and require \(P_i(u)\ne\varnothing\) for every
\(u\in A\) and each relevant hypothesis. This is stimulus applicability,
not a consequence of a set-complement formula.

A stimulus definitely distinguishes a pair if every response permitted by
either hypothesis refutes the other. In Struss's stronger ``possibly''
definition, \emph{each} hypothesis has some permitted response refuting
the other. There may also be one-sided informative stimuli, where only one
hypothesis has such a response. They must not be confused with that
each-hypothesis quantifier.
\end{hbdefinition}

\begin{hbtheorem}{thm:p16-cm-discrimination}{proved=annex:p16.cm-discrimination}{Exact relational discrimination, including the asymmetric case}
At the preceding scope, all definitely distinguishing stimuli are
\[
 D_{ij}=A\setminus\pi_U(P_i\cap P_j).
\]
The each-hypothesis possibly distinguishing stimuli, including definite ones,
are
\[
 B_{ij}=A\cap\pi_U(P_i\setminus P_j)
                \cap\pi_U(P_j\setminus P_i).
\]
The strictly partial two-sided ones are \(B_{ij}\setminus D_{ij}\).
Stimuli with some informative response for at least one hypothesis are
\[
 E_{ij}=A\cap\pi_U((P_i\setminus P_j)\cup(P_j\setminus P_i));
\]
those with no informative response are \(A\setminus E_{ij}\).
One-sided informative stimuli are \(E_{ij}\setminus B_{ij}\).
For \(P_1(u)=\{a\},P_2(u)=\{a,b\}\), the stimulus is one-sided:
\(b\) refutes hypothesis one, but no response can identify hypothesis one.
\end{hbtheorem}

This preserves the quantifier in Definition~10.16. The printed Lemma~10.3
complement expression classifies all nondefinite stimuli with at least one
exclusive response, including the one-sided case. It therefore does not
characterize the stronger each-hypothesis definition without another symmetry
condition. The distinction above resolves that mismatch rather than changing
the definition silently. For example, overlapping response sets
\(\{a,c\},\{b,c\}\) are two-sided partial; \(\{a\},\{b\}\)
are definite; identical nonempty sets provide no discrimination.

\begin{hbtheorem}{thm:p16-cm-test-sets}{proved=annex:p16.cm-test-sets}{Pair coverage and its actual hitting-set obligation}
For finitely many fixed hypotheses with the stated response semantics,
a stimulus set \(T\subseteq A\) guarantees pairwise discrimination
exactly when it hits every \(D_{ij}\), \(i<j\). Its inclusion-minimal
members are exactly the minimal hitting sets of that family. If a pair has
empty \(D_{ij}\), no such static guaranteed test set exists. If the family
is empty, the empty set is the unique inclusion-minimal answer. With
nonnegative additive test costs, a cheapest finite solution, if one exists,
has an inclusion-minimal solution of no greater cost.
\end{hbtheorem}

The construction is: compose each hypothesis, project to the declared
observables, compute pair intersections/differences, project to the stimuli,
apply admissibility restrictions, then choose a set meeting the pair
requirements. It does not need an exhaustive diagnosis model beforehand.
But explicit relations can be huge; a compact decision-diagram implementation
or search can be useful, and compilation itself may be intractable.
Selecting one witness from each nonempty pair set gives a covering set;
it does not compute a minimum-cardinality or minimum-cost set.

Static response relations also suppress hidden state and order. If several
tests can discriminate only through a persistent internal state, include
explicit state hypotheses in the product \(Hyp\times Hyp_{state}\).
Persistence through the tests is then a premise. Intrusive probing can modify
connections or components; attaching a meter can create a new circuit, so
the original \(R_i\) must not be reused without that structural revision.
Safety may exclude otherwise informative stimuli.

\section{Information-based selection and diagnosability}
\label{sec:p16-cm-information}

Relational possibility alone gives no probabilities. For the source's
information-based selection, supply finite hypothesis probabilities \(p_i\)
and a calibrated conditional response distribution \(q_i(z\mid u)\)
on a finite response alphabet.
It must agree with the permitted responses, but agreement alone does not
determine its positive weights. A test stimulus here is selected independently
of the unknown hypothesis; new observations update the probabilities.
Use \(0\log 0=0\) and omit zero-probability observations in conditional
expressions \cite[secs.~10.5.2--10.5.3]{Struss2008ModelBasedProblemSolving}.

\begin{hbtheorem}{thm:p16-cm-information}{proved=annex:p16.cm-information}{The full finite expected-information calculation}
Let \(r(z\mid u)=\sum_i p_iq_i(z\mid u)\), and for positive \(r\)
let \(p_i(z,u)=p_iq_i(z\mid u)/r(z\mid u)\). The expected reduction
in hypothesis entropy equals
\[
 H(p)-\sum_z r(z\mid u)H(p(\cdot\mid z,u))
   =H(r(\cdot\mid u))-\sum_i p_iH(q_i(\cdot\mid u)).
\]
For GDE's special approximation, each hypothesis predicts one value or
assigns a uniform distribution on \(m\ge1\) values. If the total mass of the
uniform group is \(p_{unk}\), maximizing information is equivalent to
minimizing \(\sum_z r(z)\log r(z)+p_{unk}\log m\).
The equality does not make uniformity or independence follow from an unknown
fault model, and maximizing one-step information is not a proof of optimal
total test cost.
\end{hbtheorem}

\begin{hbexample}{def:p16-cm-information-case}{Distinct predictions versus an unpredictable observation}
Take two equiprobable hypotheses. Test one deterministically returns
different binary values for them. Its observation entropy is one bit and
conditional entropy is zero, giving one bit of information. Test two
returns a fair binary outcome under \emph{both} hypotheses. Its observation
entropy is also one bit, but its conditional entropy is one bit, so the
information is zero and the posterior stays equiprobable. Thus maximizing
response variability alone is wrong. Different execution costs or risks
also require a decision criterion beyond this one-step entropy calculation.
\end{hbexample}

Probe selection fixes the stimulus and varies observable variables.
Sensor-placement and diagnosability analysis apply the same pair relation
construction to possible operating conditions: a sensor fault may show only
in a temperature range, a gear-box fault only in second gear. Fault detection
compares fault and normal models; fault discrimination compares two fault
models. In discrete-event systems, the question concerns distinguishable
observable traces within a specified delay, not a single snapshot relation.
The observables and allowed conditions are essential parts of each claim
\cite[sec.~10.5.4]{Struss2008ModelBasedProblemSolving}.

\section{Remedies, reconfiguration and process-oriented therapy}
\label{sec:p16-cm-remedies}

Diagnosis serves an action: discriminating further is unnecessary if the
same safe remedy works for all remaining cases. In a workshop, disassembly,
instrumentation and replacement can dominate the cost of reasoning.
On-board recovery may instead need only distinctions between safety classes.
The source separates replacement, switching-based reconfiguration, altered
operating regions, modified control, and structural interventions
\cite[pp.~446--449]{Struss2008ModelBasedProblemSolving}.

An integrated description includes observation, mode, action and goal
histories. Struss's simplified remedy definition assumes instantaneous,
successful actions. A consistency remedy permits the goals; an explanatory
remedy entails them. An action proposal for several mode histories must
meet its criterion for each history. In the entailment version one must
also require a consistent, applicable action model: an impossible action
must not ``achieve'' every goal by classical explosion.

\begin{hbdefinition}{def:p16-cm-reconfiguration}{State choices and what persists}
Fix a diagnostic mode assignment \(a\), a set of reconfigurable
components with exhaustive finite state domains, and persistent observations
\(O_P\subseteq O\). A target state assignment \(s_G\) is a consistency
reconfiguration when \(L\cup S\cup\{a,s_G\}\cup O_P\cup G\)
is consistent. Relative to a known actual state assignment \(s_A\), its
inclusion change set records the components given different target states.
If actual states are uncertain, common entailed states and admissible changes
must be specified rather than treating every unknown state as cost-free.
Determining \(O_P\) is itself a frame/persistence problem.
\end{hbdefinition}

Opening a breaker changes voltages/currents but need not change generator
availability. Retaining old voltages can prevent a genuine reconfiguration;
forgetting generator availability can connect consumers to an inactive source.
For a fixed hypothesis, the corresponding state-conflict clause requires at
least one currently inconsistent switch state to change. Every proposed
change must be recomposed and checked, since connecting one customer can
disconnect another. Ranking is based on change cost, not a universal
normal-versus-faulty distinction among operating states
\cite[pp.~450--453]{Struss2008ModelBasedProblemSolving}.

\begin{hbtheorem}{thm:p16-cm-remedies}{proved=annex:p16.cm-remedies}{Existence, robust outcomes and the planning boundary}
With the exhaustive state domains just stated, some complete consistency
target assignment exists for \(a\) exactly when
\(L\cup S\cup\{a\}\cup O_P\cup G\) is consistent.
The persistent observations cannot in general be omitted from this test.

For a candidate target/intervention \(d\), let \(R_i(d)\) be the
permitted outcome set under remaining hypothesis \(i\), including the
action semantics and persistent data, and let \(G_i\) be the goal-satisfying
outcome set. A consistency remedy for all hypotheses
requires \(R_i(d)\cap G_i\ne\varnothing\) for every \(i\).
A model-guaranteed remedy requires \(R_i(d)\ne\varnothing\) and
\(R_i(d)\subseteq G_i\) for every \(i\). The latter implies the
former, not conversely. For finite intervention/change sets and a cost
nondecreasing under inclusion, a cheapest feasible remedy has an
inclusion-minimal feasible remedy of no greater cost. These are goal/outcome
conditions; a valid target configuration need not be reachable by a plan
which satisfies intermediate restrictions.
\end{hbtheorem}

For example, an open-line relation which leaves the consumer voltage
unrestricted can be consistent with voltage present without guaranteeing
it. A source-connected supply condition is additional physical structure,
not something an entailment test can invent. Incomplete observation and
weak causal/behavioral models can therefore compromise remedy conclusions.
A network topology avoiding every still-suspect line may restore supply
without deciding which line is shorted, but the action sequence must also
avoid overload and maintain critical customers throughout. Activating a
generator before closing a route may be essential; target-state analysis
does not supply that ordering automatically.

\begin{hbexample}{def:p16-cm-iron-therapy}{The intervention extends the process model}
For the iron assessment, an action trigger introduces ozone or chlorine
at a specified location and enables oxidation, generating iron oxide.
The library must model this intervention and the triggered process.
The revised model must resolve the competing influences to obtain the desired
negative dissolved-iron derivative; simply having an oxidizer object is
not that inference. Under the source's continuous-quantity intervention
assumption, existing concentration does not instantly jump below its limit;
the derivative can change while the amount at the action time persists.

Therapy searches for a minimal set of permitted action triggers, not minimal
unexplained disturbance objects. A proposal to stop the pumps may reduce
iron delivery but violate the continuing water-supply goal. Several goals
may require trade-offs or several steps. Dosing in the reservoir and dosing
in the treatment tank are distinct located interventions. The chapter's
snapshot therapy formulation does not plan that multi-step history
\cite[pp.~453--454]{Struss2008ModelBasedProblemSolving}.
\end{hbexample}

\section{Design, failure analysis and software are different revision spaces}
\label{sec:p16-cm-design}

General design can change structure and is unbounded in principle.
Configuration and parametric design constrain the alternatives; a predecessor
product supplies a useful initial candidate. A configuration is a structure
and parameter specification consistent with domain constraints and configuration
goals. Such constraints often encode permissible type/parameter combinations,
not explicit models of every physical behavior. Conflict extraction can
identify design decisions requiring revision, but it does not create an
adequate initial structural hypothesis
\cite[secs.~10.2.3,10.7.1]{Struss2008ModelBasedProblemSolving}.

FMEA fixes a design and a fault/scenario set, commonly single faults with
selected multiple-fault cases. Under a consistent fault/scenario model,
if \(M_F\) is the relevant fault model and
\(C_{scen}\) the scenario, an effect \(E\) is predicted when
\(M_F\cup C_{scen}\entails E\), and possible when
\(M_F\cup C_{scen}\cup E\) is consistent. Detectability, severity and
fault probability are separate inputs to criticality assessment. Comparing
fault behavior with the normal model can produce functionally relevant
deviations, but an arbitrary numerical difference is not automatically a
safety effect \cite[sec.~10.7.2]{Struss2008ModelBasedProblemSolving}.

Software debugging violates the common artifact-diagnosis premises: a needed
computation may be missing, connections may be wrong, and all existing code
units can execute their definitions while the program fails its specification.
Goals therefore must be represented by specifications or test expectations.
The language semantics determines connections between code units. Structural
bugs cannot generally be fixed by declaring one existing line an unrestricted
faulty component. Fault catalogs cannot exhaust programmers' possible design
mistakes. Knowledge-base debugging similarly uses positive and negative
examples as goals \cite[secs.~10.7.3]{Struss2008ModelBasedProblemSolving}.

\section{Model construction is not solved by the reasoning engine}
\label{sec:p16-cm-challenges}

The chapter's concluding assessment is its 2008 historical assessment,
not a present-day claim about deployed products. It emphasizes reusable
library construction, structural/non-component faults, distributed diagnosis
and scale. Generality, adequate task resolution and compact runtime models
can pull in different directions. Automatically selecting distinctions and
compiling tailored system models from maintained fragments is substantive
work, not a consequence of giving a library a uniform syntax.

The work process matters as much as the isolated inference: required inputs
must be obtainable, output concepts must match engineers' decisions, and
testing/repair plans must reduce the costly activities and downtime.
Neither a full consistency decision nor a most-probable diagnosis by itself
supplies a reusable library, a safe test apparatus, or an executable repair
workflow \cite[sec.~10.8]{Struss2008ModelBasedProblemSolving}.

\begin{hbmethod}{def:p16-cm-consultation}{A model-based problem-solving specification}
Separate the domain library, actual structure, task and observation/goal
criterion. State the conceptual alternatives and their finite or generative
revision space. Require valid coverage of relevant environments and typed
connections before composing a candidate. Distinguish predictor soundness,
consistency completeness and conflict extraction. Record mode preferences,
unknown faults, temporal assumptions and current observation scope.
Construct tests from actual observable behavior differences, keeping
applicability, risk and hidden state explicit. For remedies, specify what
persists and whether goals are possible or guaranteed for every remaining
case; finally check reachability and intermediate restrictions with a planner.
\end{hbmethod}

\index{model-based problem solving}\index{model composition}
\index{diagnosis!process oriented}\index{diagnosis!temporal}
\index{test generation}\index{reconfiguration}\index{therapy proposal}
